\section*{Related Work}

Standardized resistance catalogues, most notably the World Health Organization's catalogue of mutations, formalize the genotype-to-phenotype mapping as a curated, rule-based lookup \cite{who2023catalogue}. Catalogues of this kind are transparent in the sense that each listed mutation carries an explicit, documented resistance association and an expert-assigned evidence grade, and independent evaluation on regional cohorts is how their diagnostic performance has been established \cite{garciamarin2024role}. They do not, however, produce a quantitative, per-isolate explanation of a prediction, cannot easily represent the combined effect of multiple co-occurring variants, and require periodic manual re-curation as new mutations are characterized.

Several mature tools instead predict resistance directly from sequence. TB-Profiler applies a curated variant library across 13 drugs \cite{phelan2019tbprofiler}, GenTB pairs random forest and wide-and-deep neural models over a comparable panel \cite{groschel2021gentb}, and DeepAMR treats co-occurrent resistance as a multi-task problem using a denoising autoencoder \cite{yang2019deepamr}. These systems set the performance context for the present work, and we compare against their reported figures in the Discussion rather than here. That comparison is complicated by the absence of a shared benchmark: reported performance for resistance prediction in this organism is strongly sensitive to which database and which isolate set is used \cite{ngo2019benchmarking}, so numbers from different studies are not directly commensurable.

Machine learning classifiers trained on genomic features can capture combinations of variants and continuous-valued statistical structure that a fixed mutation list does not, but have historically been criticized for functioning as opaque predictors that a clinician or microbiologist cannot audit. SHAP (SHapley Additive exPlanations) addresses this by attributing each individual prediction to a signed, per-feature contribution that sums to the model's output, drawing on a game-theoretic foundation for feature attribution \cite{lundberg2017unified}. Applying SHAP to a tree-based ensemble trained on genomic features makes it possible to ask, for a specific isolate, which genomic positions drove a specific resistance prediction, and to check those positions against known resistance biology. The guarantees carry conditions: the standard estimators assume features can be varied independently, an assumption genomic data satisfies only imperfectly, and we return to the consequences in the Discussion \cite{aas2021explaining,kumar2020problems}.

A growing body of work uses LLM-powered agents to assist in explainability tasks, pairing a reason-then-act loop with domain-specific tools, retrieval over external knowledge, and memory of earlier decisions. Saheed and Chukwuere apply this pattern to anomaly detection in IoT and industrial IoT networks, combining an edge-deployed agent with SHAP attributions and reporting that the agent generalizes to previously unseen attack classes where static classifiers do not \cite{saheed2026llmagent}. The components that make such a system work, namely tool use over an already-trained model, retrieval against a reference catalogue, and a record of what was asked and answered, map onto what a resistance prediction needs in order to be explained to a microbiologist. To our knowledge the pattern has not been applied to antimicrobial resistance prediction from bacterial genomes, and that gap motivates the present work.
